\chapter{Bevezetés}
\label{ch:intro}

\section{Motiváció és problémafelvetés}

Az egyetemi évek alatt számos olyan tárgyat tanulunk, amelyek megértése
jelentős absztrakciós képességet igényel, és sokszor nagyon elvontnak,
nehezen megfoghatónak tűnnek. Hallgatóként én is többször találkoztam olyan
tananyaggal, ahol a fogalmak és folyamatok elsőre nehezen voltak értelmezhetők.
Különösen igaz ez azokra a témákra, amelyekben egy rendszer működésének megértéséhez
több egymásra épülő állapotot vagy lépést kell végigkövetni.
Az ilyen absztrakt gondolatmenetek papíron gyakran nehezen követhetők,
hiszen egy folyamat vagy modell működését látni egészen más élmény,
mint pusztán olvasni róla.

Erre jó példa a Számításelmélet alapjai I--II. tárgy,
amely olyan témaköröket dolgoz fel, mint a formális nyelvek,
a grammatikák, a véges automaták, a reguláris kifejezések,
a veremautomaták vagy a Turing-gépek.
Ezek a fogalmak a hallgatók számára gyakran nehezen értelmezhetők, mivel
működésük és viselkedésük hagyományos tananyagok segítségével nem mindig
szemléltethető egyszerűen.
A tananyag elsajátításához általában előadásanyagok, prezentációk és
órai jegyzetek állnak rendelkezésünkre. Bár ezek fontos szerepet
töltenek be az oktatásban, és ezek adják a tudásunk alapját, a
mélyebb megértés sokszor a gyakorlatokra, illetve az önálló otthoni
tanulásra hárul.

Saját tapasztalatom szerint ezeknél a témáknál a megértés kulcsa
éppen az, hogy az ember minden állapotot és lépést végigkövet. Papíron
ez azonban lassú: gyakran minden állapotot külön le kell rajzolni, egy
elrontott lépés után pedig elölről kell kezdeni. Ezért az ilyen
témáknál mi, hallgatók gyakran nagyon sok időt töltünk az anyag önálló
feldolgozásával.

A szakdolgozat erre a problémára kíván megoldást nyújtani egy olyan
webalapú oktatási rendszer fejlesztésével, amely a tantárgy
legfontosabb témaköreit interaktív, vizualizációra épülő formában
dolgozza fel. A rendszer a témakörökhöz elméleti összefoglalót,
lépésenként követhető szimulációkat (például determinisztikus véges
automata, grammatikai levezetés) és önellenőrző feladatokat kínál,
így a lépésenkénti követés gyorssá és hibatűrővé válik.
% TODO: a felsorolt példákat a végén egyeztetni az elkészült modulokkal

\section{Célkitűzések}
A szakdolgozat célja egy interaktív és vizualizációs oktatási webalkalmazás
tervezése és fejlesztése, amely támogatja a hallgatókat a Számításelmélet
alapjai I--II. tantárgy témaköreinek elsajátításában.
A fejlesztés során olyan rendszer létrehozása a cél,
amely egyetlen felületen egyesíti az elméleti háttéranyagot, a vizuális
szemléltetést és az interaktív gyakorlási lehetőségeket. A hallgatók így
nem csupán olvashatnak az egyes modellekről és algoritmusokról, hanem
működésüket lépésről lépésre követhetik, és saját bemenetek segítségével
ki is próbálhatják.
A rendszer fontos célja továbbá, hogy támogassa az önálló tanulást,
valamint segítséget nyújtson a gyakorlásban és a vizsgákra való
felkészülés során.
A kitűzött célok az alábbiakban foglalhatók össze:

\begin{itemize}
    \item a Számításelmélet alapjai I--II. főbb témaköreinek feldolgozása;
    \item elméleti összefoglaló a tárgy jelöléseivel, matematikai képletekkel;
    \item saját felhasználói bemenetre működő, lépésenként követhető szimulációk;
    \item azonnali visszajelzés és önellenőrzés: kvízek, feladatok;
    \item böngészőből, telepítés nélkül elérhető felület;
    \item bővíthető, moduláris alkalmazásarchitektúra kialakítása.
\end{itemize}

\section{Hatókör és korlátok}

\subsection{Hatókör}

% TODO: a számokat és a modullistát a végén egyeztetni a kóddal
% (packages/content/src/topics, apps/web/src/modules/index.tsx)
A szakdolgozat egy webalapú, egyfelhasználós oktatási alkalmazás
tervezését és megvalósítását foglalja magában, amely a Számításelmélet
alapjai I--II. tárgy témakörét dolgozza fel Minden témakörhöz elméleti összefoglaló
tartozik, amelyet több helyen interaktív modul egészít ki.

rendszer összesen x témakört dolgoz fel, amelyek xx nagyobb csoportba sorolhatók: 
alapfogalmak, reguláris nyelvek, környezetfüggetlen nyelvek és környezetfüggő nyelvek...
 Az egyes témakörökhöz elméleti összefoglalók, valamint interaktív modulok tartoznak, 
 amelyek lehetővé teszik a modellek és algoritmusok működésének lépésről lépésre történő megfigyelését.
 
A megvalósított modulok közé tartozik többek között a DVA-szimulátor, a grammatikalevezető, a Chomsky-osztályozó, a 3-as normálforma vizualizáló, a reguláris kifejezés ellenőrző, a nyelvfelfedező, a szójáték és a kvízrendszer.

% TODO: NDVA-szimulátor elkészült? Ha igen, kerüljön a fenti listába
% TODO: ha a témakör-táblázat bekerül a fejlesztői dokumentációba, ide hivatkozás:
% A témakörök és modulok részletes áttekintését a \ref{tab:temakorok}.~táblázat tartalmazza.

A dolgozat része a fejlesztői dokumentáció, vagyis
a rendszer architektúrájának bemutatása, a tervezés
dokumentálása, az implementáció ismertetése, valamint a megvalósított
funkciók tesztelése és értékelése. Végül egy felhasználói dokumentáció
is készül, amely a rendszer használatát segíti.

\subsection{Korlátok}

A dolgozat keretében megvalósított rendszer oktatási célú
prototípusnak tekinthető. Bár alkalmas a tanulás támogatására, nem
célja egy teljes körű e-learning vagy tanulmányi rendszer kiváltása.
A rendelkezésre álló idő és fejlesztési kapacitás miatt a tantárgy teljes
 tananyagának részletes feldolgozása nem valósul meg. 
 A rendszer elsősorban azokra a témakörökre összpontosít, amelyek különösen jól támogathatók vizualizációval
  és interaktív szimulációkkal.
A Számításelmélet alapjai II. tárgy csak részben kerül feldolgozásra.
A szimulációk kisméretű, véges bemenetekre lettek tervezve.
Nincs felhasználókezelés, a haladás nem tárolódik központilag.

\subsection{Feltételezések}

A rendszer használata modern webböngésző meglétét feltételezi. A
felhasználók rendelkeznek alapvető informatikai ismeretekkel, valamint
a Számításelmélet alapjai I--II. tárgy követéséhez szükséges
előismeretekkel. Mivel a tananyag az előadásokra épül, azokat nem
helyettesíti, hanem kiegészíti. A tervezés során feltételezem, hogy a webes
technológiák és a kliensoldali vizualizációs megoldások elegendő
teljesítményt biztosítanak az interaktív modulok futtatásához.

\section{A dolgozat fő eredményei}

\begin{itemize}
    \item egy interaktív, webalapú oktatási rendszer tervezése és
          megvalósítása a Számításelmélet alapjai I--II. tárgy támogatására;
    \item olyan interaktív vizualizációs és szimulációs eszközök fejlesztése, amelyek
     lehetővé teszik a hallgatók számára a számításelméleti modellek működésének lépésről 
     lépésre történő követését;
    \item a tananyag egységes, típusos, blokkalapú tartalommodellbe
          szervezése;
    \item moduláris és bővíthető szoftverarchitektúra kialakítása;
    \item interaktív szimulációs és kiértékelő mechanizmusok fejlesztése;
    \item a rendszer tervezésének, implementációjának és tesztelésének
          dokumentálása.
\end{itemize}

A dolgozat eredménye egy olyan oktatási segédeszköz, amely a
hagyományos tananyagokat kiegészítve támogatja az önálló tanulást, a
szemléletesebb megértést és a gyakorlati tapasztalatszerzést.

\section{A dolgozat felépítése}
A dolgozat két fő részre tagolódik.
A felhasználói dokumentáció a rendszer felépítését és használatát
ismerteti a végfelhasználó szemszögéből. Ez kiterjed az egyes modulok
bemutatására és az interaktív, szimulációs eszközök használatára
példákon keresztül.

A fejlesztői dokumentáció a műszaki hátteret mutatja be. Itt kap helyet
a követelményelemzés, a tervezési döntések indoklása, az architektúra
ismertetése és az implementáció bemutatása. Ezenkívül ez a rész tér ki
a tesztelési folyamatokra, valamint a jövőbeni fejlesztési irányokra
és bővítési lehetőségekre is.
A dolgozatot egy rövid értékelés zárja.
